\section{结论}

本文提出PC-RCO，一种结合连续Gauss-Markov正则化和自适应更新率限幅的CFD先验海流观测器，用于水动力模型失配下的鲁棒AUV海流估计。主要贡献和发现如下：

\begin{enumerate}[itemsep=0pt,topsep=4pt]
    \item \textbf{模型失配鲁棒性}：在30\% CFD阻力扰动下，PC-RCO仅退化42\%，而CFD增广EKF退化77\%，鲁棒性优势在50\%失配时扩大至88个百分点。零失配下PC-RCO精度比运动学基线高3倍（RMSE$_{V_c}=0.064$ vs.\ 0.192）。

    \item \textbf{连续GM正则化}：以连续GM衰减替代二值激励门控，消除了阈值调谐需求，在失配和高噪声条件下提供57--218\%的精度改善。实验证明二值激励门控无法在反馈控制AUV中可靠区分激励水平——这一发现对本领域观测器设计具有参考意义。

    \item \textbf{在线噪声估计的自适应限幅}：单步更新上限根据实时DVL噪声估计从0.06动态调整至0.18 m/s，相比固定限幅在高噪声下RMSE降低72\%（0.86$\rightarrow$0.24），同时将无限制变体的单步跳变降低29--33倍。

    \item \textbf{多目标评估框架}：RMSE单独可能对海流观测器评估产生误导——无限制变体实现更低RMSE（0.126 vs.\ 0.243）但产生物理不可能的2.9 m/s单步估计跳变。建议将单步更新量、过冲和最终偏差作为RMSE的常规补充指标。
\end{enumerate}

未来工作将致力于：（1）以独立ADCP测量作为海流真实值的水池和海试验证；（2）将PC-RCO集成到闭环MPC中，量化有界海流估计的下游控制收益；（3）扩展为基于累积运行数据自适应更新CFD先验，以应对生物附着和结构变化导致的长期模型漂移。
